\subsection{Naturalidad de la escala y formas canónicas operatoriales}
\label{rev5:naturalidad-areal}

La longitud generada permite precisar la compatibilidad entre la
realización positiva, el operador de área y la información sectorial.
La representación marcada \(\mathcal F(K,\xi)\) conserva,
entre los datos pertinentes de \(\xi\), la resolución
\(\Omega\), sus proyectores, el retorno transportado y la sección
\(L_*\). Las coordenadas regionales generadas, la acción y los canales
constitutivos mantienen su genealogía anterior. La inversa marcada
restituye así todos los argumentos de
\eqref{rev5:longitud-generada} y \eqref{rev5:G-area}.
En particular, el lector geométrico ya construido de
\(\alpha^{\mathrm{geom}}\) determina las expresiones
\[
 L_\alpha^{\mathrm{geom}}
 =L_*(\alpha^{\mathrm{geom}})^{16}\frac{5759}{23040},
 \qquad
 G^{\mathrm{geom}}
 =\frac{(c^{\mathrm{geom}})^3
 (L_\alpha^{\mathrm{geom}})^2}{\hbar_{\rm ret}^{\mathrm{geom}}}.
\]
Los lectores de velocidad y acción se obtienen por la misma
composición con la inversa marcada. Estas expresiones conservan
los datos del retorno y la sección radial de referencia.

\begin{corollary}[Conservación por cambios de coordenatización y refinamiento heredado]
\label{rev5:conservacion}
Los cambios positivos de coordenatización que conservan el punto marcado y los
refinamientos que restituyen el registro inicial conservan
\(L_\alpha\), \(G\), \(\gamma_{\rm HMT}\) y
\(\widehat{\mathcal A}=L_\alpha^2a_\partial N\), cuando
transportan asimismo los datos sectoriales y dimensionales declarados.
El espectro normalizado
\(\widehat{\mathcal A}/L_\alpha^2=a_\partial N\), sus
multiplicidades y el estado \(\rho_A(\Delta)\) permanecen iguales.
\end{corollary}
\begin{proof}
El teorema~\ref{thm:compat-algebra-lectores} transporta cada lector
por composición con la inversa. Las mutaciones del mismo punto
reconstruyen el mismo registro. En el refinamiento heredado, sumar
las separaciones subdivididas reconstruye cada coordenada inicial.
La tensorización de la inscripción conserva \(r_\Omega\),
y sus restantes argumentos quedan fijados por las marcas. Por
sustitución se conservan la longitud, el coeficiente gravitatorio y
el operador de área. Los proyectores de ocupación conservan el
polinomio de multiplicidades \((1+z+z^2)^{36}\) y la expresión
normalizada del estado.
\end{proof}

El refinamiento heredado considerado aquí conserva el mismo espacio
de treinta y seis enlaces de frontera. Una ampliación que introduzca
un espacio físico de frontera diferente requiere entrelazadores
especificados para \(N_{90}\) y \(N_{120}\), junto con el
transporte de los datos sectoriales declarados.

Sea ahora \(P\) un politopo de rango uno de este manuscrito,
con una subdivisión orientada \(P=\bigcup_\sigma P_\sigma\)
del mismo soporte, multiplicidad uno y caras compatibles. Se fija
la misma fibra sectorial \(\mathscr H_A\) sobre todas sus
celdas. La forma de grado máximo con valores operatoriales y
dimensión de área se define por
\[
 \boldsymbol\Omega^{\mathcal A}_P
 =\widehat{\mathcal A}\otimes\Omega_P.
\]
Aquí \(\widehat{\mathcal A}\) es constante respecto de las
coordenadas diferenciales del politopo: procede del estado HMT
fijado, mientras \(\Omega_P\) describe su soporte geométrico.
El coeficiente \(\widehat{\mathcal A}\) es semidefinido positivo.
Su componente en el sector de vacío \(N=0\) es nula. Sobre un
autoespacio de ocupación \(n>0\), dividir por
\(nL_\alpha^2a_\partial\) recupera la forma canónica
\(\Omega_P\) multiplicada por la identidad de ese autoespacio.

\begin{proposition}[Aditividad areal y residuos de frontera]
\label{rev5:forma-areal}
En la trivialización sectorial común se cumplen
\begin{equation}
 \boldsymbol\Omega^{\mathcal A}_P
 =\sum_\sigma\widehat{\mathcal A}\otimes\Omega_{P_\sigma},
 \qquad
 \operatorname{Res}_F\boldsymbol\Omega^{\mathcal A}_P
 =\widehat{\mathcal A}\otimes\Omega_F.
\end{equation}
Las identidades persisten bajo subdivisiones sucesivas y residuos
iterados. Para un estado fijo \(\rho\), tomar la traza da
\(\Tr(\rho\widehat{\mathcal A})\Omega_P\) y conserva
las mismas identidades escalares.
\end{proposition}
\begin{proof}
El espacio de operadores de \(\mathscr H_A\) es finito
dimensional. En cualquier base, la afirmación es la identidad de
aditividad de las formas canónicas multiplicada por cada entrada
constante de \(\widehat{\mathcal A}\). El residuo y la traza
son lineales. La cancelación de facetas interiores conserva sus
orientaciones y la misma entrada operatorial en ambos lados.
La iteración aplica esas operaciones a cada subdivisión y a cada
bandera de caras.
\end{proof}

Si distintas celdas llevan coeficientes operatoriales regulares
\(A_\sigma\), el defecto de una pareja sobre su cara común es
\((A_\sigma|_F-A_\tau|_F)\otimes\Omega_F\).
La igualdad de restricciones operatoriales a la frontera proporciona
el pegado, bajo las condiciones de polos simples del
teorema~\ref{thm:compat-pegado}. La suma global conserva además
todas las contribuciones sobre cada divisor ambiental y el control
de polos exteriores. La normalización dimensional, la cancelación
de residuos y la cobertura del soporte tienen así criterios
separados y compatibles.

La identidad modular finita también conserva una expresión directa
en áreas físicas \(\mathcal A_m=L_\alpha^2a_\partial N_m\):
\begin{equation}
 s(\sigma)-s(\rho_0)
 =\sum_{m=90,120}\frac{m\Delta_0}{L_\alpha^2a_\partial}
   \Tr[(\sigma-\rho_0)\mathcal A_m]
  -D(\sigma\Vert\rho_0).
 \label{rev5:modular-fisica}
\end{equation}
Es \eqref{rev4:ae:ley-finita-formula} después de insertar la
longitud construida. El factor angular \(a_0\), contenido en
\(a_\partial=\gamma_{\rm HMT}a_0\), permanece íntegro.
Si se comparan secciones dimensionales con \(L'_\alpha=sL_\alpha\)
y se conserva el mismo estado y los mismos datos adimensionales,
las áreas se multiplican por \(s^2\) y sus coeficientes modulares
por \(s^{-2}\). La entropía y la entropía relativa permanecen
iguales. Esta comparación de secciones no constituye otra solución
para los mismos datos longitudinales fijados.

Para un horizonte esférico de radio \(R_h=L_\alpha\zeta_h\),
el área geométrica vale \(4\pi_{\rm HMT}L_\alpha^2\zeta_h^2\).
Su identificación con la esperanza del operador sectorial, si se
requiere para un mismo corte, exige la condición explícita
\[
 a_\partial\Tr(\rho N)=4\pi_{\rm HMT}\zeta_h^2.
\]
La geometría del horizonte y el estado de frontera deben satisfacer
esa relación con sus propias construcciones. El factor
\(\zeta_h\) caracteriza ese horizonte; la escala longitudinal
del acoplamiento permanece \(L_\alpha\).
